% Auto-generated complete chapter from Research/deep-research-report_04.md
\chapter{Benchmark salarial y trabajo remoto (Colombia)}
\label{complete:research_057}

\begin{sourcemeta}
\textbf{Fuente:} \href{file:///E:/Laboral/Research/deep-research-report_04.md}{Research/deep-research-report\_04.md}\\
\textbf{Seccion:} research\\
\textbf{Estado:} research\_snapshot\\
\textbf{Rol:} Research / estudio de mercado\\
\textbf{Lineas originales:} 89\\
\textbf{Ultima modificacion:} 2026-06-11 12:48\\
\textbf{Deduplicacion conservadora:} 0 bloques repetidos omitidos; 0 caracteres repetidos no reimpresos.
\end{sourcemeta}

\section*{Resumen de orientacion}
Este documento integra datos recientes (2025–2026) de salarios y contratación remota para perfiles 3D/Unity y roles afines. Incluye: rangos salariales por rol y región, tarifas de contratistas, costos contratistas/EOR, y detalles fiscales colombianos. Las cifras provienen de fuentes verificadas (Howdy, Glassdoor, PayScale, ERI SalaryExpert, Tecla, etc.) y reflejan situaciones reales de mercado.

\section*{Contenido completo depurado}
\addcontentsline{toc}{section}{Contenido completo depurado}




Este documento integra datos recientes (2025–2026) de salarios y contratación remota para perfiles 3D/Unity y roles afines. Incluye: rangos salariales por rol y región, tarifas de contratistas, costos contratistas/EOR, y detalles fiscales colombianos. Las cifras provienen de fuentes verificadas (Howdy, Glassdoor, PayScale, ERI SalaryExpert, Tecla, etc.) y reflejan situaciones reales de mercado.  




\subsection{Resumen de hallazgos clave}




\begin{itemize}
\item \textbf{Brecha US vs LATAM:} En EE.UU. un \textbf{Desarrollador Unity/3D en tiempo real} senior suele ganar en torno a \$100–140 mil/año (por ejemplo, Glassdoor reporta sueldo medio \textasciitilde{}\$113 K), mientras que en América Latina el promedio es mucho menor (\textasciitilde{}\$55–67 K/año según Howdy). Contratando en LATAM se ahorra \textasciitilde{}60\% comparado con EE.UU..  
\item \textbf{Europa Occidental:} En Alemania el salario promedio para \textbf{Unity/3D} es \textasciitilde{}€74 K/año, en España \textasciitilde{}€50 K. Técnicos con especialidad en Unity reciben similares (€70–75 K/año en Alemania). Países como Países Bajos o Nordics pagarían algo más (\textasciitilde{}€80–100 K), pero no hallamos datos precisos.  
\item \textbf{Roles técnicos y especializados:} Un \textbf{Artista Técnico Unity (Technical Artist)} en EE.UU. ronda \$75–80 K, y en Alemania \textasciitilde{}€70 K. Un \textbf{desarrollador WebGL/3D} en startups US promedia \textasciitilde{}\$113 K (rango muy amplio), similar a un Unity Dev. Roles XR/AR/VR en EE.UU. promedian \textasciitilde{}\$113 K (Glassdoor) – son nicho pero cotizan alto (hasta \textasciitilde{}\$150 K). Roles de \textbf{herramientas/pipeline o Python/IA} se ubican en rangos similares a desarrolladores back-end (\textasciitilde{}\$60–90 K medio en EE.UU.), aunque no hay datos específicos detallados.  
\item \textbf{Full-stack y 3D Artist (fallbacks):} Full-stack en Colombia se estima \textasciitilde{}\$30–78 K (según nivel), muy por debajo de Europa/EE.UU. Un \textbf{3D Artist/Blender} generalista ganaría menos que desarrollador (tal vez \$30–50 K en EEUU, mucho menos en LATAM). Estos roles no son prioridad principal: sirven de respaldo en caso necesario.  
\item \textbf{Contratación remota y EOR:} Cada vez más empresas usan \textbf{EOR} (Employer of Record) por temas de cumplimiento. Trabajar como contratista independiente desde Colombia conlleva retener impuestos (aprox. 28–30\% para salud/pensión) y posiblemente IVA 19\% (aunque se puede tramitar exención por “exportación de servicios” con RUT exportador). Plataformas globales (Wise, PayPal) facilitan pagos internacionales.  
\item \textbf{Metas de ingreso:} Un ingreso remoto bruto de \textbf{USD 1.5K/mes} (18K/año) es factible en 0–3 meses contratando desde LATAM (sueldo junior local o contrato medio); \textbf{USD 3K/mes} (36K/año) es meta 3–6 meses con buen portfolio (mid‑junior internacional, freelance fuerte); \textbf{USD 6K/mes} (72K/año) es difícil a corto plazo (necesita rol US/EU senior, o proyecto digital twin/IA con fuerte tracción). Tras obtener ciudadanía portuguesa (o visa EU), el rango salarial objetivo subiría, pues permitiría roles mejor pagados en Europa.
\end{itemize}




\subsection{Salarios por rol y región}




A continuación, se resumen los rangos salariales anuales (brutos) por rol y región. Los valores son aproximados y varían según experiencia, industria y empresa.  




\begin{itemize}
\item \textbf{Desarrollador 3D/Real-Time (Unity/C\#, WebGL):}  
\item \emph{EE.UU.:} base típico de \textbf{\$80K–100K} anual para perfil mid-level; seniors pueden llegar a \textbf{\$115K–150K+}. Glassdoor reporta rango \$60K–\$106K (mediana \textasciitilde{}\$80K).  
\item \emph{Alemania:} promedio \textasciitilde{}\textasciitilde{}€74K/año (rango \textasciitilde{}€50K–90K).  
\item \emph{España/Portugal:} promedio \textasciitilde{}\textasciitilde{}€50K/año (junior €35K, senior €57K). Portugal es similar a España o algo menos.  
\item \emph{LatAm (Colombia):} Unity Dev promedio \textasciitilde{}COP 108M/año (\textasciitilde{}USD \$30K); entry \textasciitilde{}COP 76M (\$21K), senior \textasciitilde{}COP 124M (\$34K). Howdy estima ingeniero LATAM \textasciitilde{}USD 55K–67K en general.  
\item \emph{Remoto global (contractor):} mercados freelance pagan típicamente +30–60\% que LATAM interno. Plataformas como Lemon indican para un Unity Senior remotos \textasciitilde{}US\$57–74/h (\textasciitilde{}\$120K–150K año full-time).  
\item \emph{Notas:} Roles Unity/WebGL suelen requerir portafolio con demos (ej. WebGL build), repositorios; TwinSight X500 es prueba fuerte para este perfil.
\end{itemize}




\begin{itemize}
\item \textbf{Artista Técnico (Technical Artist Unity):}  
\item \emph{EE.UU.:} promedio \textasciitilde{}\$75K/año (rango típicamente \$60–119K según experiencia).  
\item \emph{Alemania:} \textasciitilde{}\textasciitilde{}€70K/año (rango €49K–86K).  
\item \emph{España:} no hay datos específicos, posiblemente €40K–50K rango medio.  
\item \emph{LatAm:} menores, quizá USD 20K–40K; Howdy destaca que roles TI/Ingeniería lideran salarios, pero TA suele verse más como rol 3D/artístico. ARA (framework Python) no es factor directo en salario TA.  
\item \emph{Notas:} El portafolio de TwinSight X500 demuestra capacidad TA (optimización, UI técnica); se espera evidencia de trabajo en shaders, pipelines o herramientas para aplicar a este rol.
\end{itemize}




\begin{itemize}
\item \textbf{Desarrollador WebGL/3D (navegador interactivo):}  
\item \emph{EE.UU.:} Wellfound reporta salario medio \textasciitilde{}\textbf{\$113 K/año} para WebGL en startups (rango amplio \$38K–\$245K). Estos roles son híbridos web/3D.  
\item \emph{Europa:} se ubican en similar que Unity Dev (tal vez €60K–€100K para senior). No hay fuente específica.  
\item \emph{LatAm:} similar a Unity Dev LATAM (\textasciitilde{}USD 20–30K).  
\item \emph{Notas:} Experiencia concreta con WebGL (ej. builds en Unity o Three.js) aumenta competitividad. La vacante suele requerir demo en navegador (TwinSight X500 es ejemplo relevante).
\end{itemize}




\begin{itemize}
\item \textbf{Desarrollador XR/AR/VR:}  
\item \emph{EE.UU.:} Glassdoor estima promedio \textbf{\$112.7 K/año} (rango típico \$86K–\$149K). VR/AR es especialidad en auge.  
\item \emph{Alemania/EU:} posiblemente €60K–€100K para senior, pero datos limitados.  
\item \emph{LatAm:} faltan datos claros; probablemente mas bajos (USD 20–40K).  
\item \emph{Notas:} Cualquiera experiencia en AR/VR (por ejemplo, Unity para Meta Quest, etc.) es un plus salarial (ver Lemon: AR/VR +12–20\% premium). Sin embargo, rol XR puro es secundario al perfil principal.
\end{itemize}




\begin{itemize}
\item \textbf{Desarrollador de Herramientas/Pipeline (Python, Data, QA):}  
\item \emph{EE.UU.:} similar a desarrollador back-end: \textasciitilde{}\$60–100K dependiendo de especialidad. Python/LLM dev puede superar, pero no hay cifra concreta.  
\item \emph{Alemania:} \textasciitilde{}\$50–80K.  
\item \emph{LatAm:} \textasciitilde{}\$20–40K. Por ejemplo, desarrolladores Python/jr en Latinoamérica suelen cobrar 15–25 mil USD para juniors, y los seniors de software \textasciitilde{}60K (incluyendo mercado general).  
\item \emph{Notas:} El proyecto ARA (LangGraph) es muestra de habilidades Python/IA, útil para roles de automatización de investigación o prototipos, pero no para un rol centralizado sin portafolio de producto.
\end{itemize}




\begin{itemize}
\item \textbf{Full-Stack / Backend (fallback):}  
\item \emph{EE.UU.:} \textasciitilde{}\$80K base (rango \$60–120K+ según experiencia).  
\item \emph{LatAm:} Tecla sitúa full-stack Colombia en \textasciitilde{}USD 30–105K según nivel (junior \textasciitilde{}\$30–42K, senior \textasciitilde{}\$78–105K).  
\item \emph{Notas:} Aunque Alexander tiene conocimientos de JS/Python (ej. FitApp, AI News), este no debe ser enfoque principal. Sin embargo, fortalecer backend (p.ej. API robusta) puede abrir roles en empresas no de 3D.
\end{itemize}




\begin{itemize}
\item \textbf{3D Artist / Modelador (fallback):}  
\item \emph{EE.UU.:} un 3D generalista suele ganar \$40–60K en empresas pequeñas; en industrias como publicidad podría llegar a \$70K. No hay fuente fiable.  
\item \emph{LatAm:} mucho menor, tal vez USD 10–20K (salaries de artistas 3D freelancers en LATAM suelen ser bajos).  
\item \emph{Notas:} El retrato hiperrealista en Blender es demo de habilidades 3D, pero no es prueba de workflow de juego/pipeline. Es más bien un activo secundario de portafolio y no una ruta de carrera central.
\end{itemize}




Las \textbf{metas de ingreso} propuestas (USD 1.5K, 3K, 6K mensuales) dependen de modalidad y experiencia:




\begin{itemize}
\item \textbf{USD 1.5K/mes (\textasciitilde{}18K/año)}: viable con roles junior en LATAM (por ejemplo full-stack JR) o subcontratos locales. Se requiere CV+portafolio básico y buena red de aplicaciones freelance.  
\item \textbf{USD 3K/mes (\textasciitilde{}36K/año)}: factible en 3–6 meses apuntando a contratos remotos internacionales júnior fuerte o mid-level latino. Ejemplo: contratista remoto medio en Latinoamérica según Tarifa Howdy \textasciitilde{}USD 60–75K/año (≈\$5–6K/mes) (full cost), lo cual deja al trabajador algo menos neto.  
\item \textbf{USD 6K/mes (\textasciitilde{}72K/año)}: difícil a corto plazo. Requeriría rol norteamericano/EU de nivel medio-alto o senior (por ejemplo, \textbf{Digital Twin/Simulación}, o \textbf{Unity Senior en empresa europea}). En 12–24 meses, con portafolio sólido (TwinSight, validación académica) y tal vez ciudadanía portuguesa, podría alcanzarse en Alemania o Portugal.
\end{itemize}




\subsection{Contratación remota y modalidades}




\begin{itemize}
\item \textbf{Contratos locales vs internacionales:} Trabajar con \emph{empresa colombiana} implica salario local (bajo) y aportes de nómina obligatorios (aprox. 28.5\% del sueldo). Contratar \emph{directamente para empresa extranjera} como contratista B2B aumenta ingresos brutos (sin deducir impuestos en Colombia), pero exige manejar retenciones, IVA y seguridad social.  
\item \textbf{Employer of Record (EOR):} Plataformas (e.g. Deel, Oyster) permiten al empleador extranjero contratar legalmente en Colombia. Según Howdy, la tendencia es fuerte hacia EOR por cumplimiento. Como ventaja, el contratista recibe beneficios similares a local (EPS, pensión retenidos) pero suele sacrificar algo del ingreso neto (ya que el empleador paga \textasciitilde{}20.5\% adicional de SS). Sin embargo, EOR ofrece estabilidad y suele ser mejor opción para roles full-time remotos formales.  
\item \textbf{Tarifas de contratistas:} A nivel mundial, los seniors Unity/3D cobran \textasciitilde{}\$50–75/h (≈\$100K–150K año). En LATAM, según una guía de reclutamiento, \emph{junior} puede rondar \$25K/año y \emph{senior} \$60–75K. Por ejemplo, un desarrollador Python senior en Colombia podría negociar \textasciitilde{}\$3,000–4,000 USD al mes como contratista (antes de gastos locales). Estas tarifas deben ajustarse por la competencia y especialización.  
\item \textbf{Efecto neto (bruto a neto):} Hay que reservar \textasciitilde{}30–40\% del ingreso bruto para impuestos y SS. Por ejemplo, USD 3,000 brutos mensuales (\textasciitilde{}COP 10.5M) implicarían cotizaciones y provisiones (salud 12.5\%, pensión 16\%, ARL 0.5–6\% según riesgo) por unos COP 3–4M, dejando neto \textasciitilde{}USD 1,800–2,100. Además conviene usar medios de pago internacionales (Wise, Payoneer) para minimizar costos de cambio.
\end{itemize}




\subsection{Recomendaciones estratégicas}




\begin{itemize}
\item \textbf{Prioridad 0–3 meses:} Enfocarse en roles remotos de \textbf{Unity/WebGL Interactivo} o \textbf{Desarrollo 3D en tiempo real} nivel junior-medio. Utilizar portafolio existente (TwinSight X500, demos en WebGL) para aplicar. Ofertas en agencias LATAM (ej. cerca de \$3–5K brutos al mes) o empresas pequeñas startup. Consolidar perfil GitHub con ejemplos de optimización, shaders, interactividad 3D.  
\item \textbf{3–6 meses:} Apuntar a posiciones que valoren su especialidad “CAD-to-realtime” y UI técnico (p.ej. simulación, visualización industrial). Podría incluir Ingeniería Digital/Digital Twin (automatización, GIS, IoT). El pasaporte portugués futuro abre roles en UE; en este periodo todavía sería remoto. Expandir red en job boards especializados (p.ej. Unity Connect, StackOverflow Jobs).  
\item \textbf{12–24 meses:} Con pasaporte PT (\textasciitilde{}2027) o gracias a reagrupación familiar (Alemania), moverse a roles internos en la UE: e.g. Unity Developer sénior, Technical Artist en industria 3D (AR/VR industrial, automoción). Esto justificaría aspirar a USD 6K+/mes y beneficios completos UE (salud, pensión, protección laboral). Mientras tanto, prepararse para entrevistas internacionales (idiomas, certificados de Blender/Unity).
\end{itemize}




En resumen, la ruta más coherente es enfocarse en \textbf{Desarrollador 3D/Unity/WebGL} (perfil híbrido dev+TA) como objetivo principal, usando TwinSight como carta de presentación. Full-stack y 3D generalista quedan como respaldo secundario. Las oportunidades en IA/automatización (proyecto ARA) pueden servir para roles de tooling interno, pero no deben desviar el foco principal 3D.




\subsection{Próxima sección}




\begin{mdcode}
Bloque de codigo: markdown
04_eu_portugal_germany_mobility_strategy.md
\end{mdcode}




Para la estrategia de movilidad (ciudadanía portuguesa, matrimonios, visados) recomiendo activar \textbf{Deep Research}. Es necesario consultar fuentes oficiales (SEF Portugal, Gobierno Alemán, web de visados UE) para obtener requisitos, tiempos y costos actualizados (2024–2026). No usaremos modo agente aún; el perfil de estudios requiere más análisis documental que recolección masiva.



